\documentclass{beamer}
\usetheme{Madrid}

% Encoding and fonts for pdfLaTeX
\usepackage[utf8]{inputenc}
\usepackage[T1]{fontenc}
\usepackage{lmodern}
\usepackage{hyperref}

\title{Prompt Engineering}
\author{Mehmet Ali Akyol, PhD}
\institute{Ankara Medipol University}
\date{\today}

\begin{document}

\begin{frame}
    \titlepage
    \begin{center}
        \vspace{1cm}
        \textbf{Week 4: Prompt Patterns \& Frameworks}
    \end{center}
\end{frame}

\section{Overview}

\begin{frame}
    \frametitle{Where We Are}
    \begin{itemize}
        \item Week 1: Prompt fundamentals and the power of structured instructions.
        \item Week 2: The mechanics of modern language models.
        \item Week 3: Principles and anti-patterns for effective prompt design.
        \item \textbf{This week:} Codifying best practice with reusable prompt patterns and frameworks.
    \end{itemize}
\end{frame}

\begin{frame}
    \frametitle{Weekly Learning Objectives}
    \begin{itemize}
        \item Differentiate between atomic prompt patterns and broader frameworks.
        \item Combine patterns to address complex, multi-step tasks.
        \item Select frameworks based on task type, risk level, and time available.
        \item Translate frameworks into templates that teammates can reuse.
        \item Evaluate tradeoffs between prescriptive templates and flexible prompting.
    \end{itemize}
\end{frame}

\begin{frame}
    \frametitle{Key Concepts}
    \begin{itemize}
        \item Pattern vs. framework vs. playbook
        \item Pattern stacks (role + format + reasoning)
        \item Prompt lifecycle frameworks (PLAN, DRAFT, REVIEW)
        \item Task diagnostics (clarity, constraints, context, checkpoints)
        \item Governance patterns (guardrails, refusal handling)
    \end{itemize}
\end{frame}

\section{Pattern Library}

\begin{frame}
    \frametitle{Atomic Patterns Recap}
    \begin{itemize}
        \item \textbf{Role}: Persona, audience, tone setters.
        \item \textbf{Format}: Tables, JSON, bullet, code block.
        \item \textbf{Reasoning}: Step-by-step, chain-of-thought, critique.
        \item \textbf{Context}: Retrieved excerpts, facts, style guides.
        \item \textbf{Constraints}: Word limits, banned topics, sourcing rules.
    \end{itemize}
\end{frame}

\begin{frame}
    \frametitle{Pattern Stacking}
    \begin{itemize}
        \item Combine patterns intentionally: role + format + reasoning \textrightarrow{} reliable outputs.
        \item Use delimiters to isolate instructions, data, and examples.
        \item Add validation requests:\ ``List assumptions before final answer."\ or\ ``Return JSON conforming to schema.''
        \item Introduce fallback logic:\ ``If information missing, respond with: insufficient context.''
    \end{itemize}
\end{frame}

\begin{frame}
    \frametitle{Adaptation Patterns}
    \begin{itemize}
        \item \textbf{Surface shift}: Change tone or reading level without altering content.
        \item \textbf{Format shift}: Convert text to tables, checklists, timelines.
        \item \textbf{Perspective shift}: Rewrite from stakeholder viewpoints.
        \item \textbf{Language shift}: Translate while preserving style constraints.
    \end{itemize}
\end{frame}

\section{Frameworks}

\begin{frame}
    \frametitle{Framework Examples}
    \begin{itemize}
        \item \textbf{CO-STAR}: Context, Objective, Style, Tone, Audience, Response.
        \item \textbf{RACI Prompt Canvas}: Role, Action, Constraints, Input.
        \item \textbf{IDEA Loop}: Identify task, Draft prompt, Evaluate output, Adjust.
        \item \textbf{SMART Prompt}: Specific, Measurable, Achievable, Relevant, Time-bound outputs.
    \end{itemize}
\end{frame}

\begin{frame}
    \frametitle{Choosing a Framework}
    \begin{itemize}
        \item \textbf{Low stakes}: Lightweight templates (role + format) suffice.
        \item \textbf{Moderate stakes}: Add evaluation and guardrails (IDEA Loop, SMART Prompt).
        \item \textbf{High stakes}: Multi-stage prompts with critique, verification, and escalation rules.
        \item Consider team adoption: Can colleagues follow it quickly?
    \end{itemize}
\end{frame}

\begin{frame}
    \frametitle{Framework to Template}
    \begin{itemize}
        \item Start with fields the requester fills (objective, audience, constraints).
        \item Encode defaults for tone, format, delimiters to ensure consistency.
        \item Provide example completions highlighting best practice.
        \item Offer optional add-ons (e.g., risk checks, summary level).
    \end{itemize}
\end{frame}

\section{Deep Dive}

\begin{frame}[fragile]
    \frametitle{Deep Dive: CO-STAR Framework}
    \begin{itemize}
        \item \textbf{C}ontext: \textit{Summarize the following user feedback for a product manager.}
        \item \textbf{O}bjective: \textit{Extract the core feature request and the user's primary pain point.}
        \item \textbf{S}tyle: \textit{Use clear, professional language.}
        \item \textbf{T}one: \textit{Neutral and objective.}
        \item \textbf{A}udience: \textit{A busy product manager who needs the key takeaway in seconds.}
        \item \textbf{R}esponse: \textit{Format as a JSON object with keys "feature\_request" and "pain\_point".}
    \end{itemize}
\end{frame}

\begin{frame}[fragile]
    \frametitle{Deep Dive: RACI Prompt Canvas}
	\begin{block}{Origin}
		Adapted from the classic RACI project management matrix (Responsible, Accountable, Consulted, Informed), which has been used for decades to clarify roles in complex projects.
	\end{block}
    \begin{itemize}
        \item \textbf{R}ole: \textit{You are a helpful, expert-level programming assistant.}
        \item \textbf{A}ction: \textit{Refactor the provided Python code to be more idiomatic and efficient.}
        \item \textbf{C}onstraints: \textit{Use list comprehensions where appropriate and ensure the code adheres to PEP 8 standards.}
        \item \textbf{I}nput: \textit{The code to be refactored is provided below in the text block.}
    \end{itemize}
\end{frame}

\begin{frame}
    \frametitle{Deep Dive: The IDEA Loop}
	\begin{block}{Origin}
		A modern variant of classic continuous improvement cycles like PDCA (Plan-Do-Check-Act) or John Boyd's OODA Loop (Observe-Orient-Decide-Act), tailored for the iterative nature of prompt engineering.
	\end{block}
    \begin{itemize}
        \item \textbf{I}dentify: The goal is to get a concise, 3-bullet summary of a long report.
        \item \textbf{D}raft: The first prompt is simple: "Summarize this report."
        \item \textbf{E}valuate: The output is a long paragraph, not the desired bulleted list.
        \item \textbf{A}djust: The prompt is refined: "Summarize this report in exactly three bullet points."
    \end{itemize}
\end{frame}

\begin{frame}
    \frametitle{Deep Dive: SMART Prompt Outputs}
	\begin{block}{Origin}
		Directly borrows from the widely-used SMART criteria for effective goal-setting, first appearing in a 1981 article by George T. Doran for defining project objectives.
	\end{block}
    \begin{itemize}
        \item \textbf{S}pecific: \textit{Generate 5 multiple-choice questions about photosynthesis.}
        \item \textbf{M}easurable: \textit{Each question must have exactly 4 options: 1 correct and 3 incorrect.}
        \item \textbf{A}chievable: \textit{The questions should be suitable for a 9th-grade biology class.}
        \item \textbf{R}elevant: \textit{The questions must cover the topics of chlorophyll and cellular respiration.}
        \item \textbf{T}ime-bound: \textit{This is a single-turn request; generate all 5 questions now.}
    \end{itemize}
\end{frame}

\begin{frame}
    \frametitle{Advanced Pattern: Meta-Prompting}
    \begin{block}{Concept}
        Using a prompt to have an LLM generate, critique, or refine another prompt.
    \end{block}
    \begin{itemize}
        \item \textbf{Generator Prompt:} "Create a prompt for a junior developer to document a Python function."
        \item \textbf{Critique Prompt:} "Review this prompt. Is it clear? Does it specify the output format? Suggest one improvement."
        \item \textbf{Refinement Prompt:} "Given the following critique, revise the original prompt to be more effective."
        \item \textbf{Use Case:} Automating the creation and maintenance of a high-quality prompt library.
    \end{itemize}
\end{frame}

\begin{frame}
    \frametitle{Advanced Pattern: Tool Use \& Function Calling}
    \begin{block}{Concept}
        Empowering an LLM to use external tools by providing a manifest of available functions and their schemas.
    \end{block}
    \begin{itemize}
        \item The model doesn't execute code; it generates a JSON object specifying the tool to call and the arguments to use.
        \item \textbf{Example Flow:}
        \begin{enumerate}
            \item User asks: "What's the weather in Istanbul?"
            \item Model returns: \texttt{\{"tool": "get\_weather", "city": "Istanbul"\}}
            \item Your code executes the function and passes the result back to the model.
            \item Model generates a natural language response based on the tool's output.
        \end{enumerate}
    \end{itemize}
\end{frame}

\begin{frame}
    \frametitle{Advanced Pattern: Dynamic Few-Shot Selection}
    \begin{block}{Concept}
        Instead of using static, hard-coded examples, programmatically select the most relevant examples from a library based on the user's query.
    \end{block}
    \begin{itemize}
        \item \textbf{How it works:}
        \begin{enumerate}
            \item Store a large library of high-quality prompt/completion pairs.
            \item When a user query comes in, use a similarity search (e.g., vector embeddings) to find the most similar examples in your library.
            \item Dynamically insert these top K examples into the prompt before sending it to the LLM.
        \end{enumerate}
        \item \textbf{Benefit:} Provides highly relevant, in-context learning for the model on-the-fly.
    \end{itemize}
\end{frame}

\begin{frame}
    \frametitle{From Frameworks to Playbooks}
    \begin{block}{Concept}
        A playbook is an organizational asset that goes beyond a single prompt, detailing a complete, multi-step workflow for a complex business process.
    \end{block}
    \begin{itemize}
        \item \textbf{A Playbook Contains:}
        \begin{itemize}
            \item The business goal and success metrics.
            \item A sequence of prompts (a prompt chain).
            \item Human-in-the-loop review and approval steps.
            \item Guardrails and escalation paths for failures.
            \item Examples of good and bad final outputs.
        \end{itemize}
        \item \textbf{Goal:} To make AI-driven processes reliable, repeatable, and scalable across a team.
    \end{itemize}
\end{frame}

\section{Worked Examples}

\begin{frame}
    \frametitle{Policy Brief Template}
    \begin{itemize}
        \item Context: national AI policy update, 1-page summary.
        \item Framework: CO-STAR + critique loop.
        \item Output: Executive summary, key impacts, stakeholder risks, action items.
        \item Guardrail: require citation markers linked to provided context IDs.
    \end{itemize}
\end{frame}

\begin{frame}
    \frametitle{Customer Support Escalation}
    \begin{itemize}
        \item Objective: classify tickets and draft replies.
        \item Pattern stack: role + format + reasoning + guardrail.
        \item Steps: detect sentiment, map to policy, produce compliant response, list disclaimers.
        \item Evaluation: JSON report with confidence fields and suggested follow-up.
    \end{itemize}
\end{frame}

\begin{frame}
    \frametitle{Research Synthesis Pattern}
    \begin{itemize}
        \item Inputs: annotated bibliography, research question, stakeholder.
        \item Framework: PLAN (Purpose, Lens, Audience, Next steps).
        \item Output: 3-section memo with citations, gaps, recommended experiments.
        \item Reflection: require identification of uncertain claims for human review.
    \end{itemize}
\end{frame}

% --- Activities Section Removed ---

% --- Assignment Section Removed ---

\section{Wrap-up}

\begin{frame}
    \frametitle{Key Takeaways}
    \begin{itemize}
        \item Patterns are building blocks; frameworks package them for reuse.
        \item Match framework complexity to task stakes and resources.
        \item Templates accelerate teams when paired with clear guardrails.
        \item Iterate: test, critique, and refine before broad rollout.
    \end{itemize}
\end{frame}

\begin{frame}{Questions and Discussion}
    \begin{center}
    {\Huge Thank You!}
    \vspace{1cm}

        \textbf{Dr.\ Mehmet Ali Akyol}\\
        \texttt{mehmetali.akyol@gmail.com}
    \vspace{1cm}

    {\large Questions?}
    \end{center}
\end{frame}

\end{document}
